% !TEX program = xelatex
\documentclass[12pt]{ctexart}

\usepackage[a4paper,margin=2.15cm,top=2.55cm,bottom=2.25cm]{geometry}
\usepackage{amsmath,amssymb}
\usepackage{enumitem}
\usepackage{array}
\usepackage{fancyhdr}

\setlist[enumerate]{itemsep=9pt,topsep=5pt,leftmargin=2.2em}
\linespread{1.13}

\pagestyle{fancy}
\fancyhf{}
\lhead{\small 高等代数 I 期末复习进阶卷 · A 卷}
\rhead{\small 第 \thepage 页}
\setlength{\headheight}{14pt}
\renewcommand{\headrulewidth}{0.4pt}

\newcommand{\rnk}{\operatorname{r}}
\newcommand{\blank}[1][2.6cm]{\underline{\hspace{#1}}}
\newcommand{\bigsection}[2]{%
  \vspace{7pt}%
  {\noindent\bfseries\large #1}\hfill{\normalsize（本大题共 #2）}\par
  \vspace{2pt}\hrule height0.4pt\vspace{8pt}}
\newcommand{\optgrid}[4]{%
  \par\smallskip\noindent
  \begin{tabular}{@{}p{0.46\linewidth}p{0.46\linewidth}@{}}
  (A)~#1 & (B)~#2 \\[3pt]
  (C)~#3 & (D)~#4
  \end{tabular}}
\newcommand{\workarea}[1]{\par\vspace{0.35em}\noindent\rule{\linewidth}{0.25pt}\par\vspace{#1}}

\begin{document}

\begin{center}
  {\Large\bfseries 高等代数 I 期末复习进阶卷（A 卷）}\\[4pt]
  {\small 教材：北京大学数学系《高等代数》（第 6 版）为主\\
  范围：行列式、线性方程组、向量组线性相关性、矩阵与逆矩阵、初等变换与初等矩阵、矩阵与向量组的秩}\\[8pt]
\end{center}

\noindent
\begin{tabular}{@{}p{4.6cm}p{4.6cm}p{4.6cm}@{}}
姓名：\blank[2.8cm] & 学号：\blank[2.8cm] & 班级：\blank[2.8cm]
\end{tabular}

\vspace{6pt}
\noindent\textbf{满分 100 分\qquad 考试时间 120 分钟\qquad 闭卷}

\vspace{8pt}
\noindent
\renewcommand{\arraystretch}{1.4}
\begin{tabular}{|c|c|c|c|c|c|}
\hline
题号 & 一 & 二 & 三 & 四 & 总分 \\\hline
满分 & 30 & 20 & 10 & 40 & 100 \\\hline
得分 & & & & & \\\hline
\end{tabular}

\vspace{12pt}

\bigsection{一、单项选择题}{10 小题，每题 3 分，共 30 分}

\begin{enumerate}[label=\arabic*.]
\item 在 $5$ 阶行列式的展开式中，项 $a_{12}a_{23}a_{35}a_{41}a_{54}$ 所带的符号为（\qquad）
  \optgrid{正号 $+$}{负号 $-$}{$0$}{无法确定}

\item 设 $A$ 为 $n$ 阶方阵，且 $A^{2}-3A+2E=O$，则下列结论正确的是（\qquad）
  \optgrid{$A$ 不一定可逆}{$A^{-1}=\dfrac{3E-A}{2}$}{$A^{-1}=A-3E$}{$|A|=0$}

\item 设 $A$ 为 $m\times n$ 矩阵，$B$ 为 $n\times s$ 矩阵，则必有（\qquad）
  \optgrid{$\rnk(AB)\geq \rnk(A)+\rnk(B)-n$}{$\rnk(AB)=\rnk(A)+\rnk(B)-n$}{$\rnk(AB)\geq \rnk(A)+\rnk(B)-m$}{$\rnk(AB)=\min\{\rnk(A),\rnk(B)\}$}

\item 设 $A$ 为 $n$ 阶矩阵，$\rnk(A)=n-1$，$A^{*}$ 为其伴随矩阵，则 $\rnk(A^{*})=$（\qquad）
  \optgrid{$0$}{$1$}{$n-1$}{$n$}

\item 对矩阵 $A$ 施行一次初等列变换，相当于（\qquad）
  \optgrid{左乘相应的初等矩阵}{右乘相应的初等矩阵}{两侧同乘相应的初等矩阵}{取转置后再作行变换}

\item 设向量组 $\alpha_1,\dots,\alpha_s$ 线性无关，且 $\beta$ 不能由它们线性表示，则（\qquad）
  \optgrid{$\alpha_1,\dots,\alpha_s,\beta$ 线性无关}{$\alpha_1,\dots,\alpha_s,\beta$ 线性相关}{$\beta=0$}{$s$ 必等于向量空间维数}

\item 非齐次线性方程组 $Ax=b$（$A$ 为 $m\times n$ 矩阵）有无穷多解的充要条件是（\qquad）
  \optgrid{$\rnk(A)<\rnk(A,b)$}{$\rnk(A)=n$}{$\rnk(A,b)=n$}{$\rnk(A)=\rnk(A,b)<n$}

\item 设 $A_{jk}$ 为 $n$ 阶行列式中元素 $a_{jk}$ 的代数余子式。当 $i\neq j$ 时，
  $\displaystyle\sum_{k=1}^{n}a_{ik}A_{jk}=$（\qquad）
  \optgrid{$|A|$}{$0$}{$1$}{$-|A|$}

\item $n$ 元齐次线性方程组 $Ax=0$ 的系数矩阵秩为 $n-2$，则它的基础解系含（\qquad）个向量。
  \optgrid{$0$}{$1$}{$2$}{$n-2$}

\item $n$ 阶矩阵 $A$ 可以表示成若干个初等矩阵的乘积的充要条件是（\qquad）
  \optgrid{$A$ 可逆}{$|A|=0$}{$\rnk(A)<n$}{$A$ 为对角矩阵}
\end{enumerate}

\bigsection{二、填空题}{10 小题，每题 2 分，共 20 分}

\begin{enumerate}[label=\arabic*.]
\item 排列 $6\,2\,5\,1\,4\,3$ 的逆序数为 $\blank[2cm]$。
\item $4$ 阶行列式
  $\begin{vmatrix}
  a&b&b&b\\ b&a&b&b\\ b&b&a&b\\ b&b&b&a
  \end{vmatrix}=\blank[4cm]$。
\item $\begin{vmatrix}1&1&1\\ x&y&z\\ x^{2}&y^{2}&z^{2}\end{vmatrix}=\blank[4cm]$。
\item 若 $P,Q$ 均为 $2$ 阶方阵，$|P|=2,\ |Q|=-3$，则
  $\begin{vmatrix}P&O\\ R&Q\end{vmatrix}=\blank[2.5cm]$。
\item 设 $A$ 为 $4$ 阶方阵且 $|A|=2$，则 $|A^{*}|=\blank[2.2cm]$。
\item 设 $A$ 可逆且 $AB=O$，则 $B=\blank[2cm]$。
\item 设 $A$ 为 $5\times 7$ 矩阵，$\rnk(A)=3$，则齐次方程组 $Ax=0$ 的解空间维数为 $\blank[2cm]$。
\item 若 $A$ 为 $n$ 阶可逆矩阵且 $A^{2}-2A=O$，则 $A=\blank[2cm]$。
\item 若 $\alpha_1,\alpha_2,\alpha_3,\alpha_4$ 线性无关，则向量组
  $\alpha_1+\alpha_2,\alpha_2+\alpha_3,\alpha_3+\alpha_4,\alpha_4$ 的秩为 $\blank[2cm]$。
\item 若非齐次线性方程组 $Ax=b$ 有解，且 $\rnk(A)=\rnk(A,b)=r<n$，则通解中自由未知量个数为 $\blank[2cm]$。
\end{enumerate}

\bigsection{三、证明题}{2 小题，每题 5 分，共 10 分}

\begin{enumerate}[label=\arabic*.]
\item 设 $A,B$ 为 $n$ 阶方阵，$AB=O$，且 $A+B$ 可逆。证明：
  \[
  \rnk(A)+\rnk(B)=n.
  \]
\workarea{4.2cm}

\item 设向量组 $\alpha_1,\dots,\alpha_s$ 线性无关，向量 $\beta$ 不能由 $\alpha_1,\dots,\alpha_s$ 线性表示。证明：
  $\alpha_1,\dots,\alpha_s,\beta$ 线性无关。
\workarea{4.2cm}
\end{enumerate}

\clearpage
\bigsection{四、计算题}{4 小题，每题 10 分，共 40 分。要求写出主要步骤}

\begin{enumerate}[label=\arabic*.]
\item 用拉普拉斯定理或分块思想计算行列式
  \[
  D=\begin{vmatrix}
  a&b&0&0\\
  c&d&0&0\\
  1&2&a&b\\
  3&4&c&d
  \end{vmatrix}.
  \]
\workarea{5.2cm}

\item $a$ 取何值时，线性方程组
  \[
  \begin{cases}
  a x_1+x_2+x_3=1,\\
  x_1+a x_2+x_3=1,\\
  x_1+x_2+a x_3=a^{2}
  \end{cases}
  \]
  无解、有唯一解、有无穷多解？在有无穷多解时求出通解。
\workarea{6.0cm}

\clearpage
\item 用初等行变换求矩阵
  \[
  A=\begin{pmatrix}
  1&2&1\\
  2&5&2\\
  1&0&3
  \end{pmatrix}
  \]
  的逆矩阵 $A^{-1}$。
\workarea{6.0cm}

\item 设向量组
  \[
  \alpha_1=(1,0,0,1),\quad \alpha_2=(0,1,0,1),\quad
  \alpha_3=(0,0,1,1),
  \]
  \[
  \alpha_4=(1,1,1,3),\quad \alpha_5=(2,1,0,3).
  \]
  求该向量组的秩与一个极大线性无关组，并把其余向量用该极大无关组线性表示。
\workarea{5.5cm}
\end{enumerate}

\begin{center}\small 试题部分结束，参考答案见后。\end{center}

\clearpage
\begin{center}
  {\Large\bfseries 参考答案与评分标准（A 卷）}\\[2pt]
  {\small 高等代数 I 期末复习进阶卷}
\end{center}
\vspace{6pt}\hrule\vspace{12pt}

\noindent\textbf{一、单项选择题}（每题 3 分）
\begin{center}
\renewcommand{\arraystretch}{1.3}
\begin{tabular}{|c|c|c|c|c|c|c|c|c|c|c|}
\hline
题号 & 1&2&3&4&5&6&7&8&9&10\\\hline
答案 & A&B&A&B&B&A&D&B&C&A\\\hline
\end{tabular}
\end{center}

\noindent\small 简析：1. 列指标排列 $(2,3,5,1,4)$ 的逆序数为 $4$，故为正号。2. 由
$A(3E-A)=2E$ 得 $A$ 可逆，$A^{-1}=(3E-A)/2$。3. 为 Sylvester 秩不等式。
4. $\rnk(A)=n-1$ 时伴随矩阵秩为 $1$。5. 初等列变换对应右乘初等矩阵。7. 有解且自由未知量个数为正。8. 代数余子式错位相乘之和为 $0$。\normalsize

\vspace{8pt}
\noindent\textbf{二、填空题}（每题 2 分）
\begin{enumerate}[label=\arabic*.,itemsep=3pt]
\item $10$。
\item $(a+3b)(a-b)^3$。
\item $(y-x)(z-x)(z-y)$。
\item $-6$。
\item $8$。
\item $O$。
\item $4$。
\item $2E$。
\item $4$。
\item $n-r$。
\end{enumerate}

\vspace{6pt}
\noindent\textbf{三、证明题}（每题 5 分）
\begin{enumerate}[label=\textbf{\arabic*.},itemsep=10pt]
\item 因 $AB=O$，$B$ 的列向量都属于齐次方程组 $Ax=0$ 的解空间，故
  $\rnk(B)\leq n-\rnk(A)$，即 $\rnk(A)+\rnk(B)\leq n$。又
  $A+B$ 可逆，所以
  \[
  n=\rnk(A+B)\leq \rnk(A)+\rnk(B).
  \]
  两式合并得 $\rnk(A)+\rnk(B)=n$。

\item 设
  $k_1\alpha_1+\cdots+k_s\alpha_s+k\beta=0$。若 $k\neq0$，则
  \[
  \beta=-\frac{k_1}{k}\alpha_1-\cdots-\frac{k_s}{k}\alpha_s,
  \]
  与 $\beta$ 不能由 $\alpha_1,\dots,\alpha_s$ 线性表示矛盾。故 $k=0$。
  于是 $k_1\alpha_1+\cdots+k_s\alpha_s=0$，由原向量组线性无关得
  $k_1=\cdots=k_s=0$。所以扩充后的向量组线性无关。
\end{enumerate}

\vspace{6pt}
\noindent\textbf{四、计算题}（每题 10 分）
\begin{enumerate}[label=\textbf{\arabic*.},itemsep=11pt]
\item 该行列式为分块下三角行列式：
  \[
  D=\begin{vmatrix}a&b\\c&d\end{vmatrix}
    \begin{vmatrix}a&b\\c&d\end{vmatrix}
  =(ad-bc)^2.
  \]

\item 系数行列式
  \[
  \Delta=\begin{vmatrix}a&1&1\\1&a&1\\1&1&a\end{vmatrix}
  =(a-1)^2(a+2).
  \]
  当 $a\neq1,-2$ 时，$\Delta\neq0$，方程组有唯一解。\\
  当 $a=-2$ 时，三式左端相加为 $0$，右端相加为 $1+1+4=6\neq0$，故无解。\\
  当 $a=1$ 时，方程组化为 $x_1+x_2+x_3=1$，有无穷多解：
  \[
  (x_1,x_2,x_3)=(1,0,0)+s(-1,1,0)+t(-1,0,1),\quad s,t\in\mathbb{R}.
  \]

\item 对 $(A\,|\,E)$ 作初等行变换可得
  \[
  A^{-1}=
  \begin{pmatrix}
  \frac{15}{2}&-3&-\frac12\\
  -2&1&0\\
  -\frac52&1&\frac12
  \end{pmatrix}
  =\frac12\begin{pmatrix}
  15&-6&-1\\
  -4&2&0\\
  -5&2&1
  \end{pmatrix}.
  \]
  评分时要求写出主要消元步骤，结果占 $3$ 分。

\item 以 $\alpha_1,\dots,\alpha_5$ 为列作矩阵
  \[
  M=\begin{pmatrix}
  1&0&0&1&2\\
  0&1&0&1&1\\
  0&0&1&1&0\\
  1&1&1&3&3
  \end{pmatrix}.
  \]
  第 $4$ 行等于前三行之和，且前三列已含 $3$ 阶单位子矩阵，故 $\rnk(M)=3$。
  可取 $\{\alpha_1,\alpha_2,\alpha_3\}$ 为极大线性无关组，并且
  \[
  \alpha_4=\alpha_1+\alpha_2+\alpha_3,\qquad
  \alpha_5=2\alpha_1+\alpha_2.
  \]
\end{enumerate}

\end{document}
